\documentclass[11pt]{article}
\usepackage{mathblatt}


\blattfuss{Quadratische}{Lösungen}
\begin{document}
\noindent{\large\bfseries Quadratische \textperiodcentered{} Lösungen}\par\medskip
\begin{pfloesung}{}
\lz{1.}{$A(3|2)$}{}{}
\end{pfloesung}
\begin{pfloesung}{}
\lz{2.}{$9$}{$g(4) = 8 + 1$}{}
\end{pfloesung}
\begin{pfloesung}{}
\lz{3.}{$x^2 + 8x + 16$}{}{}
\end{pfloesung}
\begin{pfloesung}{}
\lz{4.}{$4$}{$x$}{}
\end{pfloesung}
\begin{pfloesung}{}
\lz{5.}{$S(2|3)$}{$x + 1 = -x + 5$, $x = 2$}{}
\end{pfloesung}
\begin{pfloesung}{}
\lz{6.}{$8$}{}{}
\end{pfloesung}
\begin{pfloesung}{}
\lz{7.}{$x_1 = 6$, $x_2 = -6$}{}{}
\end{pfloesung}
\begin{pfloesung}{}
\lz{8.}{$18$}{$3 \cdot 6$}{}
\end{pfloesung}
\begin{pfloesung}{}
\lz{9.}{$x^2 + 4x$}{}{}
\end{pfloesung}
\begin{pfloesung}{}
\lz{10.}{$4$}{Scheitel bei $x$}{}
\end{pfloesung}
\begin{pfloesung}{}
\lz{11.}{Normalform}{}{}
\end{pfloesung}
\begin{pfloesung}{}
\lz{12.}{$S(2|5)$}{}{}
\end{pfloesung}
\begin{pfloesung}{}
\lz{13.}{Schnittpunkt $(0|5)$}{$f(0) = 5$}{}
\end{pfloesung}
\begin{pfloesung}{}
\lz{14.}{$S(-4|2)$}{}{}
\end{pfloesung}
\begin{pfloesung}{}
\lz{15.}{$S(2|{-3})$}{}{}
\end{pfloesung}
\begin{pfloesung}{}
\lz{16.}{$S(-2|3)$}{}{}
\end{pfloesung}
\begin{pfloesung}{{\small\color{mbgrau}P10 2017 · 5c}}
\lz{17.}{S($-$3; $-$2)}{}{}
\end{pfloesung}
\begin{pfloesung}{{\small\color{mbgrau}P10 2018 · 5b}}
\lz{18.}{S(2; $-$2)}{}{}
\end{pfloesung}
\begin{pfloesung}{{\small\color{mbgrau}P10 2020 · 3a}}
\lz{19.}{S($-$2; $-$4)}{}{}
\end{pfloesung}
\begin{pfloesung}{{\small\color{mbgrau}P10 2022 · 3b}}
\lz{20.}{S(2|$-$4)}{Scheitelpunktform p(x) = (x $-$ d)² + e mit d = 2, e = $-$4}{}
\end{pfloesung}
\begin{pfloesung}{{\small\color{mbgrau}P10 2021 · 1d}}
\lz{21.}{S($-$2; $-$1)}{}{}
\end{pfloesung}
\begin{pfloesung}{}
\lz{22.}{$(x - 5)^2 + 2$}{$f(x)$}{}
\end{pfloesung}
\begin{pfloesung}{}
\lz{23.}{$S(4|3)$}{$f(x) = x^2 - 8x + 16 - 16 + 19 = (x - 4)^2 + 3$}{}
\end{pfloesung}
\begin{pfloesung}{}
\lz{24.}{$(x - 2)^2 - 3$}{$S(2|{-3})$, nach oben geöffnet: $f(x)$}{}
\end{pfloesung}
\begin{pfloesung}{{\small\color{mbgrau}P10 2016 · 1g}}
\lz{25.}{y = (x $-$ 1)² + 3}{}{}
\end{pfloesung}
\begin{pfloesung}{}
\lz{26.}{$(x + 4)^2 + 2$, Scheitel $S(-4|2)$}{$g(x)$}{}
\end{pfloesung}
\begin{pfloesung}{}
\lz{27.}{$(x - 1)^2 - 2$, Scheitel $S(1|{-2})$}{$g(x)$}{}
\end{pfloesung}
\begin{pfloesung}{}
\lz{28.}{$-(x - 5)^2 - 1$, Scheitel $S(5|{-1})$}{$g(x)$}{}
\end{pfloesung}
\begin{pfloesung}{{\small\color{mbgrau}P10 2015 · 1i}}
\lz{29.}{S(2; 0)}{}{}
\end{pfloesung}
\begin{pfloesung}{{\small\color{mbgrau}P10 2018 · 5c}}
\lz{30.}{p(x) = (x $-$ 2)² $-$ 2}{x² $-$ 4x + 4 $-$ 4 + 2}{}
\end{pfloesung}
\begin{pfloesung}{{\small\color{mbgrau}P10 2020 · 3d}}
\lz{31.}{Skizze: Normalparabel mit Scheitel (2; $-$4); y = (x $-$ 2)² $-$ 4}{Scheitel ($-$2 | $-$4) $\Rightarrow$ (2 | $-$4)}{}
\end{pfloesung}
\begin{pfloesung}{{\small\color{mbgrau}P10 2017 · 5e}}
\lz{32.}{q: y = $-$(x + 3)²}{nach Verschiebung y = (x + 3)²}{}
\end{pfloesung}
\begin{pfloesung}{{\small\color{mbgrau}P10 2025 · 5b}}
\lz{33.}{S(3|$-$2); p(x) = (x $-$ 3)² $-$ 2}{x² $-$ 6x + 9 $-$ 2 (quadratische Ergänzung)}{}
\end{pfloesung}
\begin{pfloesung}{}
\lz{34.}{$5$}{$S(3|{-4})$; $f(x) = (x - 3)^2 - 4$; Probe: $(0 - 3)^2 - 4$; wA}{}
\end{pfloesung}
\begin{pfloesung}{}
\lz{35.}{$(x - 4)^2$ mit $S(4|0)$}{z; B. $f(x)$}{}
\end{pfloesung}
\begin{pfloesung}{{\small\color{mbgrau}P10 2023 · 4b}}
\lz{36.}{S($-$1|6); p(x) = $-$(x + 1)² + 6}{Scheitelpunktform p(x) = $-$(x $-$ d)² + e mit d = $-$1, e = 6}{}
\end{pfloesung}
\begin{pfloesung}{{\small\color{mbgrau}P10 2018 · 5d}}
\lz{37.}{Skizze: Normalparabel mit Scheitel (0; e), e > 0; z. B. q(x) = x² + 1; z. B. q$'$(x) = $-$x² $-$ 1}{e > 0, nach oben geöffnet}{}
\end{pfloesung}
\begin{pfloesung}{}
\lz{38.}{$-12$}{$p = -1$, $q$}{}
\end{pfloesung}
\begin{pfloesung}{}
\lz{39.}{$x + 4$}{$x^2 - 2x$}{}
\end{pfloesung}
\begin{pfloesung}{}
\lz{40.}{$x_1 = 3$, $x_2 = -1$}{$0 = (x - 1)^2 - 4$, $4 = (x - 1)^2$, $x - 1 = 2$ oder $x - 1 = -2$}{}
\end{pfloesung}
\begin{pfloesung}{}
\lz{41.}{$x_1 = -2$, $x_2 = -4$}{$x_{1,2} = -3 \pm \sqrt{9 - 8} = -3 \pm 1$}{}
\end{pfloesung}
\begin{pfloesung}{}
\lz{42.}{$x_1 = 4$, $x_2 = -2$}{$0 = x^2 - 2x - 8$; $x = 1 \pm \sqrt{1 + 8} = 1 \pm 3$}{}
\end{pfloesung}
\begin{pfloesung}{}
\lz{43.}{$x_1 = 3$, $x_2 = -2$}{$0 = 3x^2 - 3x - 18$ $\mid : 3$: $0 = x^2 - x - 6$; $x = 0{,}5 \pm \sqrt{0{,}25 + 6} = 0{,}5 \pm 2{,}5$}{}
\end{pfloesung}
\begin{pfloesung}{{\small\color{mbgrau}P10 2020 · 3e}}
\lz{44.}{x$_1$ = $-$1; x$_2$ = $-$3}{2x² + 8x + 6 = 0 $\Rightarrow$ x² + 4x + 3 = 0; p-q-Formel $\Rightarrow$ x = $-$2 ± 1}{}
\end{pfloesung}
\begin{pfloesung}{}
\lz{45.}{$\approx 0{,}27$}{$x = 2 \pm \sqrt{4 - 1} = 2 \pm \sqrt{3}$; $x_1 \approx 3{,}73$, $x_2$}{}
\end{pfloesung}
\begin{pfloesung}{}
\lz{46.}{$-6$}{$0 = x^2 + 12x + 36$; $x = -6 \pm \sqrt{36 - 36} = -6 \pm 0$; genau eine Nullstelle $x$}{}
\end{pfloesung}
\begin{pfloesung}{}
\lz{47.}{keine Nullstelle, denn aus $-3$ gibt es keine Wurzel}{$0 = x^2 - 4x + 7$; $x = 2 \pm \sqrt{4 - 7}$; Wert unter der Wurzel: $4 - 7 = -3$}{}
\end{pfloesung}
\begin{pfloesung}{{\small\color{mbgrau}P10 2025 · 5c}}
\lz{48.}{N1(1,59|0), N2(4,41|0), x = 3 ± √2}{Diskriminante 9 $-$ 7 = 2; √2 $\approx$ 1,41}{}
\end{pfloesung}
\begin{pfloesung}{{\small\color{mbgrau}P10 2017 · 5d}}
\lz{49.}{(x + 3)² $-$ 2 = x² + 6x + 9 $-$ 2 = x² + 6x + 7; x$_1$ = $-$3 + √2 $\approx$ $-$1,59, x$_2$ = $-$3 $-$ √2 $\approx$ $-$4,41}{(x + 3)² = 2 $\Rightarrow$ x = $-$3 ± √2}{}
\end{pfloesung}
\begin{pfloesung}{}
\lz{50.}{Formel}{Normalform}{}
\end{pfloesung}
\begin{pfloesung}{}
\lz{51.}{Wurzelziehen}{reinquadratisch}{}
\end{pfloesung}
\begin{pfloesung}{}
\lz{52.}{$S_1(5|6)$, $S_2(1|{-2})$}{$x^2 - 4x + 1 = 2x - 4$, $x^2 - 6x + 5 = 0$; $x = 3 \pm \sqrt{9 - 5} = 3 \pm 2$; $x_1 = 5$, $x_2 = 1$; $g(5) = 6$, $g(1) = -2$}{}
\end{pfloesung}
\begin{pfloesung}{}
\lz{53.}{$2$}{$x^2 + 1 = 2x$; $x^2 - 2x + 1 = 0$; $x = 1 \pm \sqrt{1 - 1}$, also $x = 1$; $g(1)$}{}
\end{pfloesung}
\begin{pfloesung}{}
\lz{54.}{$-3$, negativ}{$x^2 + 4 = 2x$; $x^2 - 2x + 4 = 0$; unter der Wurzel $1 - 4$}{}
\end{pfloesung}
\begin{pfloesung}{}
\lz{55.}{$x_1 = 3$, $x_2 = -2$}{$x^2 = x + 6$; $x^2 - x - 6 = 0$; $x_{1,2} = 0{,}5 \pm \sqrt{0{,}25 + 6} = 0{,}5 \pm 2{,}5$}{}
\end{pfloesung}
\begin{pfloesung}{}
\lz{56.}{$x_1 = 3$, $x_2 = 2$}{$(x - 2)^2 - 1 = x - 3$, $x^2 - 4x + 3 = x - 3$, $x^2 - 5x + 6 = 0$; $x = 2{,}5 \pm \sqrt{6{,}25 - 6} = 2{,}5 \pm 0{,}5$}{}
\end{pfloesung}
\begin{pfloesung}{}
\lz{57.}{$x_1 = 3$, $x_2 = -2$}{$x^2 + x - 4 = 2x + 2$, $x^2 - x - 6 = 0$; $x = 0{,}5 \pm \sqrt{0{,}25 + 6} = 0{,}5 \pm 2{,}5$}{}
\end{pfloesung}
\begin{pfloesung}{{\small\color{mbgrau}P10 2021 · 2c}}
\lz{58.}{S$_1$($-$1; $-$2); S$_2$(2; 7)}{x² + 2x $-$ 1 = 3x + 1 $\Rightarrow$ x² $-$ x $-$ 2 = 0; x$_1$ = $-$1, x$_2$ = 2; y = 3x + 1 $\Rightarrow$ y$_1$ = $-$2, y$_2$ = 7}{}
\end{pfloesung}
\begin{pfloesung}{{\small\color{mbgrau}P10 2022 · 3c}}
\lz{59.}{S$_1$($-$1|5), S$_2$(3|$-$3)}{x² $-$ 2x $-$ 3 = 0; x = 1 ± 2}{}
\end{pfloesung}
\begin{pfloesung}{{\small\color{mbgrau}P10 2024 · 3d}}
\lz{60.}{S1($-$1|$-$3), S2(5|21)}{x² $-$ 4x $-$ 5 = 0; x = 2 ± 3}{}
\end{pfloesung}
\begin{pfloesung}{}
\lz{61.}{Gerade}{}{}
\end{pfloesung}
\begin{pfloesung}{}
\lz{62.}{Parabel}{}{}
\end{pfloesung}
\begin{pfloesung}{{\small\color{mbgrau}P10 2026 FOR · 1e}}
\lz{63.}{erste Tabelle (y = 12, 3, 0, 3, 12)}{3 · 2² = 12}{}
\end{pfloesung}
\begin{pfloesung}{{\small\color{mbgrau}P10 2015 · 4b}}
\lz{64.}{h(t) = $-$3t² + 2400; h(0) = 2400 und Parabel nach unten geöffnet $\Rightarrow$ $-$3t² + 2400}{h(0) = 2400; h(20) = $-$3 · 20² + 2400 = 1200}{}
\end{pfloesung}
\begin{pfloesung}{}
\lz{65.}{$3$}{$f(x) = -x^2 + 3$: nach unten geöffnet, $f(0)$}{}
\end{pfloesung}
\begin{pfloesung}{}
\lz{66.}{$2$}{Tabelle A: $0{,}5 \cdot 2^2$}{}
\end{pfloesung}
\begin{pfloesung}{}
\lz{67.}{$4$}{$4$; $1$; $0$; $1$}{}
\end{pfloesung}
\begin{pfloesung}{}
\lz{68.}{$7x^2$}{$a = 7$, denn $f(1) = a \cdot 1^2 = 7$; $f(x)$}{}
\end{pfloesung}
\begin{pfloesung}{}
\lz{69.}{$36$}{$f(-6) = (-6)^2$}{}
\end{pfloesung}
\begin{pfloesung}{}
\lz{70.}{$-5$}{$f(-3) = 9 - 12 - 2$}{}
\end{pfloesung}
\begin{pfloesung}{{\small\color{mbgrau}P10 2020 · 3b}}
\lz{71.}{D($-$1,5; $-$2,5)}{}{}
\end{pfloesung}
\begin{pfloesung}{{\small\color{mbgrau}P10 2018 · 5a}}
\lz{72.}{wahr; falsch}{p($-$1) = 7; p(0) = 2 $\Rightarrow$ Schnittpunkt (0 | 2)}{}
\end{pfloesung}
\begin{pfloesung}{}
\lz{73.}{ja, $P$ liegt auf beiden Graphen}{$f(4) = 16 - 5 = 11$, $g(4) = 8 + 3 = 11$}{}
\end{pfloesung}
\begin{pfloesung}{{\small\color{mbgrau}P10 2024 · 3c}}
\lz{74.}{Nein, p($-$5) = 25 $-$ 4 = 21 $\neq$ 20}{($-$5)² = 25}{}
\end{pfloesung}
\begin{pfloesung}{{\small\color{mbgrau}P10 2020 · 3c}}
\lz{75.}{y = 12}{(2 + 2)² $-$ 4}{}
\end{pfloesung}
\begin{pfloesung}{{\small\color{mbgrau}P10 2014 · 7a}}
\lz{76.}{ja, S($-$3; $-$9) ist Schnittpunkt}{p($-$3) = $-$($-$3)² = $-$9; g($-$3) = 2 · ($-$3) $-$ 3 = $-$9}{}
\end{pfloesung}
\begin{pfloesung}{}
\lz{77.}{Scheitel $S(2|1)$, dann $1$ nach rechts $1$ nach oben, $2$ nach rechts $4$ nach oben, nach links genauso: $(1|2)$, $(3|2)$, $(0|5)$, $(4|5)$}{}{}
\end{pfloesung}
\begin{pfloesung}{}
\lz{78.}{Bogen durch $(-2|16)$, $(-1|4)$, $(0|0)$, $(1|4)$, $(2|16)$}{}{}
\end{pfloesung}
\begin{pfloesung}{{\small\color{mbgrau}P10 2024 · 3b}}
\lz{79.}{Parabel mit Scheitel (0|$-$4) durch (±2|0) und (±3|5); S(0|$-$4)}{Normalparabel um 4 nach unten verschoben; Nullstellen ±2}{}
\end{pfloesung}
\begin{pfloesung}{}
\lz{80.}{$c$ mit $c$ kleiner als $-2$, denn der Scheitel $S(0|{-2})$ ist der tiefste Punkt}{z; B. $g(x) = -3$; jede waagerechte Gerade $g(x)$}{}
\end{pfloesung}
\begin{pfloesung}{}
\lz{81.}{kein gemeinsamer Punkt: Scheitel $(1|2)$ und $(1|{-4})$ übereinander, gleiche Öffnung}{$g$ ist $f$ um $6$ nach unten verschoben}{}
\end{pfloesung}
\begin{pfloesung}{{\small\color{mbgrau}P10 2014 · 7b}}
\lz{82.}{z. B. f(x) = 1}{Scheitel S(0 | 0), Parabel nach unten geöffnet}{}
\end{pfloesung}
\begin{pfloesung}{}
\lz{83.}{Scheitel über der $x$-Achse, nach oben geöffnet}{Scheitel $S(2|3)$: keine Nullstelle}{}
\end{pfloesung}
\begin{pfloesung}{}
\lz{84.}{rechts positiv}{zwei Lösungen}{}
\end{pfloesung}
\begin{pfloesung}{}
\lz{85.}{rechts steht null}{gilt}{}
\end{pfloesung}
\begin{pfloesung}{}
\lz{86.}{$x_1 = 9$, $x_2 = -9$}{}{}
\end{pfloesung}
\begin{pfloesung}{}
\lz{87.}{$x_1 = 2$, $x_2 = 5$}{$x - 2 = 0$ oder $x - 5 = 0$}{}
\end{pfloesung}
\begin{pfloesung}{}
\lz{88.}{$25$ : ja}{$(-7 + 2)^2 = (-5)^2$; wA}{}
\end{pfloesung}
\begin{pfloesung}{}
\lz{89.}{$0$ : ja}{$(-5 + 5) \cdot (-5 - 1) = 0 \cdot (-6)$; wA}{}
\end{pfloesung}
\begin{pfloesung}{}
\lz{90.}{$0$ : ja}{$(-2)^2 - 3 \cdot (-2) - 10 = 4 + 6 - 10$; wA}{}
\end{pfloesung}
\begin{pfloesung}{}
\lz{91.}{$-14$}{$x = -7$, denn $(-7) \cdot (-7 + 9) = (-7) \cdot 2$}{}
\end{pfloesung}
\begin{pfloesung}{{\small\color{mbgrau}P10 2025 · 1h}}
\lz{92.}{x = $-$2}{($-$2) · 3 = $-$6}{}
\end{pfloesung}
\begin{pfloesung}{}
\lz{93.}{$x_1 = 2$, $x_2 = -3$}{$x^2 + x + 1 = 7$, $x^2 + x - 6 = 0$; $x = -0{,}5 \pm \sqrt{0{,}25 + 6} = -0{,}5 \pm 2{,}5$}{}
\end{pfloesung}
\begin{pfloesung}{}
\lz{94.}{$x_1 = 6$, $x_2 = -2$}{$-x^2 + 4x + 2 = -10$, $-x^2 + 4x + 12 = 0$ $\mid \cdot (-1)$: $x^2 - 4x - 12 = 0$; $x = 2 \pm \sqrt{4 + 12} = 2 \pm 4$}{}
\end{pfloesung}
\begin{pfloesung}{}
\lz{95.}{$x_1 = 5$, $x_2 = -4$}{$2x^2 - 2x - 40 = 0 \mid : 2$; $x^2 - x - 20 = 0$; $x_{1,2} = 0{,}5 \pm \sqrt{0{,}25 + 20} = 0{,}5 \pm 4{,}5$}{}
\end{pfloesung}
\begin{pfloesung}{{\small\color{mbgrau}P10 2023 · 4c}}
\lz{96.}{x$_1$ = $-$5, x$_2$ = 3}{x² + 2x $-$ 15 = 0; x = $-$1 ± 4}{}
\end{pfloesung}
\begin{pfloesung}{\textbf{Prüfstein} \textperiodcentered{} P10 2026 FOR · 5}
\lz{c)}{S(2|$-$2)}{Scheitelpunktform p(x) = (x $-$ d)² + e mit d = 2, e = $-$2}{}
\lz{d)}{genau ein gemeinsamer Punkt S(2|$-$2), gleicher Scheitel, p nach oben, q nach unten geöffnet $\Rightarrow$ Berührung nur im Scheitel}{Scheitel beider Parabeln (2|$-$2)}{}
\end{pfloesung}
\end{document}